


\subsection{4.1, Determinants}

\begin{problems}

\problem % Jankowski Spring 2019 quiz 6
True or false: If $A$ is a $3 \times 3$ matrix satisfying
\[Ae_1 = e_2, \qquad Ae_2 = e_2 + e_3, \qquad Ae_3 = e_2-e_3,\] then
$\det(A) = 0$.

\begin{solution}
True.  We see
    \[A = \mat{Ae_1 Ae_2 Ae_3} = \mat{0 0 0; 1 1 1; 0 1 -1}
      \xrightarrow[]{R_2 = R_2 + R_3} \mat{0 0 0; 1 2 0; 0 1 -1}.\] By
    doing just one row replacement (which does not change the
    determinant), we have obtained a triangular matrix with
    determinant $(0)(2)(-1)=0$, so $\det(A) = 0$.
\end{solution} 

\problem % Created for ILA
True or false: If $A$ is an $n \times n$ matrix and the equation $Ax = b$ has
more than one solution for some $b$ in ${\bf R}^n$, then $\det(A) = 0$.

\begin{solution}
  True.  Since $Ax=b$ has infinitely many solutions for some $b$, we know $Ax=0$
  has infinitely many solutions.  Therefore, $A$ is not invertible and so
  $\det(A)=0$.
\end{solution} 

\problem % Created for ILA
True or false: If $A$ and $B$ are $n \times n$ matrices and $\det(AB) \neq 0$,
then $A$ and $B$ must both be invertible.

\begin{solution}
  True: $\det(AB) = \det(A)\det(B)$ and the quantity is nonzero by assumption,
  so both $\det(A)$ and $\det(B)$ must be nonzero, hence $A$ and $B$ are both
  invertible.
\end{solution} 

\problem % Created for ILA
True or false: If $A$ and $B$ are $n \times n$ matrices and $\det(A)=0$, then
$AB$ is not invertible.

\begin{solution}
  True: $\det(AB) = \det(A) \det(B) = 0 \det(B) = 0$, so $AB$ cannot be
  invertible.
\end{solution} 



\problem % Dan's notes
Use row operations to find the determinant of the following matrix:
\[A = \mat{4 -3 0 1; 0 2 12 4; 0 1 -1 1; 0 2 0 1}.\]

\begin{solution} We row-reduce $A$ until we get a triangular matrix.
\begin{align*}
  \mat{4 -3 0 1; 0 2 12 4; 0 1 -1 1; 0 2 0 1}
  & \xrightarrow[]{R_2= R_2/2} \mat{4 -3 0 1; 0 1 6 2; 0 1 -1 1; 0 2 0 1}
  \xrightarrow[R_4 = R_4-2R_2]{R_3 = R_3-R_2} %
    \mat{4 -3 0 1; 0 1 6 2; 0 0 -7 -1; 0 0 -12 -3} \\
  & \xrightarrow[]{R_4 = R_4 - (12/7)R_3}
    \mat{4 -3 0 1; 0 1 6 2; 0 0 -7 -1; 0 0 0 -9/7}
\end{align*}
Let's call the final matrix $B$, which is upper-triangular (so $\det(B)$ is the
product of its diagonal entries).  To obtain $B$ from $A$, we did a row scale of
$1/2$, followed by a bunch of row replacements that don't affect determinants.
Therefore,
\[\det(A)(1/2) = \det(B), \qquad \det(A)= 2 \det(B) = 2 \cdot 4(1)(-7)(-9/7) = 72.\]
\end{solution}

\problem % Dan's notes
Find the value of $h$ that makes the determinant of the following matrix equal to $0$.
\[\mat{1 2 3; 1 0 1; 2 2 h}.\]

\begin{solution}
The determinant is $0$ if and only if the matrix is not invertible, so we row-reduce:
\[
  \mat{1 2 3; 1 0 1; 2 2 h}
  \xrightarrow[R_3 = R_3 - 2R_1]{R_2 = R_2 - R_1}
    \mat{1 2 3; 0 -2 -2; 0 -2 h-6}
  \xrightarrow[]{R_3 = R_3 - R_2}
    \mat{1 2 3; 0 -2 -2; 0 0 h-4}.
\]
This matrix is invertible unless it has fewer than $3$ pivots.  Therefore, the
determinant of the matrix is zero if and only if $h-4=0$, so $h=4$.
\end{solution}

\problem % Dan's notes
If the matrix on the left has determinant $5$, what is the determinant of the
matrix on the right?
\[\mat{a b c; d e f; g h i} \qquad
\mat{g h i; d e f; a-d b-e c-f}.\]

\begin{solution}
We track the row operations required to obtain the final matrix from the original:
\[\mat{a b c; d e f; g h i}
  \xrightarrow[\textup{new det is }-5]{R_1 \leftrightarrow R_3} \mat{g h i; d e f; a b c}
  \xrightarrow[\textup{det is still } -5]{R_3=R_3-R_2} \mat{g h i; d e f; a-d b-e c-f}.\]
\end{solution}
The first operation was a row swap which switched the determinant from $5$ to $-5$.
The second operation was a row-replacement which did not affect the determinant.
Therefore, the determinant of the final matrix is $-5$.



\problem % Dan's notes
Suppose $T$ is a linear transformation with the property that $T \circ T = T$.
What are the possibilities for the determinant of the standard matrix for $T$?

\begin{solution}
Let $A$ be the standard matrix for $T$, so $T(x) = Ax$.  Note
\[T(x) = (T \circ T)(x) = T(T(x)) = T(Ax)=A(Ax) = A^2 x,\] so $A = A^2$.  Using
this and the multiplicative property of determinants, we get
\[\det(A) = \det(A^2) = \det(A) \det(A) = (\det \ A)^2.\]
For $x = \det(A)$, the above says $x = x^2$, so $0 = x^2-x = x(x-1)$.  Therefore
$x=0$ or $x=1$, so $\det(A)$ is $0$ or $1$.

Both $0$ and $1$ are possible for the determinant of $A$.  For example, the
transformations $T(x) = \mat{0 0; 0 0}x$ and $T(x) = \mat{1 0; 0 1}x$ both
satisfy $T \circ T = T$.
%Matrices with the property $A^2 = A$ are called \textit{idempotent matrices}.  
\end{solution}


\problem % Dan's notes
Suppose $T$ is a linear transformation with the property that $T \circ T$ is the
identity transformation.  What is the determinant of the standard matrix for
$T$?

\begin{solution}
  Let $A$ be the standard matrix for $T$.  The same argument as the previous
  exercise shows that the standard matrix for $T \circ T$ is $A^2$.  But we were
  told $T \circ T$ is the identity transformation, so we also know that its
  standard matrix is $I$, therefore $I = A^2$.  Now:
  \[\det(I) = \det(A^2) \ \implies \  1 = \det(A)\det(A) \ \implies \  1 = (\det \ A)^2
    \ \implies \ \det(A) = \pm 1.\] %
  Both values for $\det(A)$ are possible in practice.  For example, the matrices
  $\mat{1 0; 0 1}$ and $\mat{1 0; 0 -1}$ have determinants $1$ and $-1$
  respectively, and both satisfy $A^2 = I$.
  %This kind of matrix is very interesting, because the fact
  %that $A^2 = I$ means that $A$ is its own inverse!
\end{solution}


\problem % Jankowski Spring 2018 e3
Let $A$ be the $3 \times 3$ matrix satisfying $Ae_1= 3e_2$, $Ae_2= e_3$, and
$Ae_3=e_1$.  Find $\det(A)$.

\begin{solution}
  $A=\mat{0 0 1; 3 0 0; 0 1 0}$.  You can compute $\det(A)=1$ or just do two row
  swaps to get a diagonal matrix, so that $\det(A)=(-1)^2(3)=3$.
\end{solution}

\problem  %made for ILA
%
Suppose $A$ is an invertible $2 \times 2$ matrix and $\det(A) = 2$.
Find $\det( (3A)^{-1} )$.

\begin{solution}
  By properties of the determinant,
  $\det( (3A)^{-1} ) = \dfrac{1}{\det(3A)}$.  Since $3A$ is obtained
  from $A$ by multiplying \textit{each row} of $A$ by $3$, we need to
  multiply the determinant by $3$ twice.
  \[\det(3A) = 3^2 \det(A) = 9 (2) = 18.\]
  Therefore,
  \[\det((3A)^{-1}) = \dfrac{1}{18}.\]
\end{solution}
          

\problem %2.2 problem in new way, which was created for ILA
A problem from section 2.2, revisited: \newline Find all
values of $c$ and $d$ so that
$\textup{Span}\left\{\vec{1 c}, \vec{d, 2}\right\}$ is not
${\bf R}^2$.

\begin{solution}
  Two vectors in ${\bf R}^2$ will fail to span ${\bf R}^2$ precisely when
  they are linearly dependent, so we solve:
  \[\det\mat{1 d; c 2} = 0, \qquad 2 - cd = 0, \qquad cd = 2.\]
  Therefore, $d = 2/c$ where $c$ is any nonzero number (the fact that
  $cd=2$ means neither $c$ or $d$ can be zero).
\end{solution}


\problem % Created for ILA
Is there a $2 \times 2$ matrix (whose entries are real numbers) with the property
that $A^2 = -I$?  If so, give an example.  If not, justify why there is no such
matrix.  

\begin{solution}
  Yes.  Take $A$ to be the matrix for rotation by $90^\circ$ counterclockwise.
  Then $A^2$ is rotation by $2 \cdot 90 ^\circ = 180^\circ$, which is the same
  as scalar multiplication by $-1$.  In other words, $A^2 = -I$.  We can verify
  this algebraically:
  \[\mat{0 -1; 1 0} \mat{0 -1; 1 0} = \mat{-1 0; 0 -1}.\]
\end{solution}

\problem % Created for ILA
Is there a $3 \times 3$ matrix (whose entries are real numbers) with the
property that $A^2 = -I$?  If so, give an example.  If not, justify why there is
no such matrix.  

\begin{solution}
  No.  If $A^2 = -I$ then \[\det(A^2) = \det(-I) = \det\mat{-1 0 0; 0 -1 0; 0 0 -1} = -1.\]
  But this would mean
  \[\det(A^2) = -1, \quad (\det A)^2 = -1, \qquad \det(A) = \pm i.\]
  This is impossible, since the determinant of every real matrix is real.
  \newline \newline It is natural to ask: why doesn't the above argument also
  show that we can never have $A^2 = -I$ in the $2 \times 2$ case?  Note that
  the $2 \times 2$ matrix $-I$ has determinant equal to $1$ (not $-1$) since
  \[\det\mat{-1 0; 0 -1} = (-1)^2 = 1,\] so we do not obtain an impossible
  equality for $\det(A)$ in the $2 \times 2$ case, and in fact we can found an
  actual example of a $2 \times 2$ matrix $A$ so that $A^2 = -I$.
\end{solution}



\problem % Created for ILA
Suppose $A$ is a square matrix and that $A^k$ is invertible for some $k$.  Must
it be true that $A$ is invertible?  Justify your answer.

\begin{solution}
  Yes, $A$ must be invertible.  We use the fact that a square matrix is
  invertible if and only if its determinant is nonzero.  Since $A^k$ is
  invertible for some $k$, this means $\det(A^k) \neq 0$.  From the properties
  of determinants we know $\det(A^k) = (\det \ A)^k$, so the previous sentence
  means $(\det A)^k \neq 0$, thus $\det(A) \neq 0$.  Therefore, $A$ is
  invertible.
\end{solution}
\end{problems}\setcounter{problem}{0}

\newpage

\subsection{4.2, Cofactor expansions}

\begin{problems}
\problem % Dan's notes
  True or false: The cofactor expansion of an $n \times n$ matrix along the
  first row gives the negative of the cofactor expansion along the second row.
  
  \begin{solution}
    False.  The cofactor expansion gives the same determinant, no matter what
    row (or column) you use.
  \end{solution}

\problem % Created for ILA
Use a cofactor expansion to find $\det(A)$ for 
$A =  \mat{5 1 -3 1; 0 2 3 4; 0 0 1 2; 5 -5 2 3}$.

\begin{solution}
We expand along the third row, which has only two nonzero entries.
\begin{align*}
  \det(A) & = 0 C_{31} + 0 C_{32} + 1 C_{33} + 2 C_{34} \\
          & = 1(-1)^{3+3} \det\mat{5 1 1; 0 2 4; 5 -5 3} +
            2(-1)^{3+4} \det \mat{5 1 -3; 0 2 3; 5 -5 2} \\
          & = 5(6+20)-1(0-20)+1(0-10) \quad + \quad
            (-2)\left(5(4+15)-1(0-15) +(-3)(0-10)\right) \\
          & = 140 - 2(140) = -140.
\end{align*}
\end{solution}

\problem % Created for ILA
As a general rule, which method for computing $n \times n$ determinants is
faster: row-reduction, or the cofactor expansion?  Justify your answer.

\begin{solution}
  Row-reduction is generally much faster, particularly if the matrix is large.
  Cofactor expansions scale very poorly as $n$ increases.  For example, in order
  to compute the determinant of a $5 \times 5$ matrix using cofactors, we need
  to compute $5$ different $4 \times 4$ determinants.  Each of these
  $4 \times 4$ determinants requires us to compute $4$ different $3 \times 3$
  determinants, and each of the $3 \times 3$ determinants requires us to compute
  $3$ different $2 \times 2$ determinants.  In essence, the $5 \times 5$
  cofactor expansion would involve $5! = 120$ computations of $2 \times 2$
  determinants.  
\end{solution}

\problem % Dan's notes
Find the determinant of the following matrix using one of the cofactor expansion formulas.
\[\mat{1 0 -2; 3 1 -2; -5 0 9}.\]

\begin{solution}
We expand along the second column, since it has only one nonzero entry.
\[\det\mat{1 0 -2; 3 1 -2; -5 0 9} = 0C_{21} + 1 C_{22} + 0C_{32}
  = 1 (-1)^{2+2} \det\mat{1 -2; -5 9} = 1(9-10) = -1.
  \]
\end{solution}

\problem % Created for ILA
Find all real values of $h$ (if there are any) so that the following set of
vectors is a basis for ${\bf R}^{3}$.
\[\left\{\vec{1 -7 1}, \vec{h 4 -1}, \vec{3 h 2}\right\}.\]
\begin{solution}
  The three vectors form a basis for ${\bf R}^{3}$ if and only if they are
  linearly independent, so put them as columns of a $3 \times 3$ matrix $A$.
  Since $A$ is square, the vectors are linearly independent if and only
  if the matrix is invertible, which is when $\det(A) \neq 0$.  Now:
  \[\det(A) = 1(8+h) - h(-14-h) + 3(7-4) = 8+h+14h+h^2 + 9
    = h^2+15h+17.\]
  This determinant is zero when
  \[h = \dfrac{-15 \pm \sqrt{15^2-4(17)}}{2} = \dfrac{-15 \pm \sqrt{157}}{2}.\]
  As established above, the vectors are linearly independent when the determinant
  is nonzero, which is all real numbers $h$ {\bf except} %
  $\dfrac{-15 + \sqrt{157}}{2}$ and $\dfrac{-15 - \sqrt{157}}{2}$.
\end{solution}

\problem % Dan's notes
Find the determinant of the following matrix using one of the cofactor expansion formulas.
\[\mat{4 0 -2; -7 1 -2; -5 -1 9}.\]
\begin{solution}
We use the first row.  The determinant of the matrix is
\[
  4\Big[1(9)-(-1)(-2)\Big] - 0 + (-2)\Big[(-7)(-1)-(-5)(1)\Big]
  = 4(9-2) - 2(7+5) = 28-24 = 4.
\]
\end{solution}

\problem %Jankowski Spring 2018 Exam 3, also Spring 2018 sample midterm 3
  Find $\det(A^3)$ if $A= \mat{1 -3 4 2; 0 0 -2 0; 0 1 2 3; 2 0 -1 20}$.  

  \begin{solution}
    Using the cofactor expansion along the second row, we find
    \[\det(A)=-2(-1)^5\det\mat{1 -3 2; 0 1 3; 2 0 20}
      = 2\Big(20+3(-6)+2(-2)\Big)=2(20-18-4)=-4,\] so
    $\det(A^3)=(-4)^3= -64$.
  \end{solution}

\problem % Created for ILA
Use the cofactor expansion of the determinant to find $\det(A)$ for
\[A = \mat{-1 3 4 0; 4 1 3 0; 1 1 4 5; -1 -3 -4 0}.\]
\begin{solution}
We use the expansion along the fourth column.
\begin{align*}
  \det(A)
  & = 0C_{14} + 0C_{24} + 5 C_{34} + 0 C_{44} \\
  & = 5 (-1)^{3+4} \det\mat{-1 3 4; 4 1 3; -1 -3 -4} \\
  & = -5\Big[-1(-4+9)-3(-16+3)+4(-12+1)\Big] \\
  & = -5(-5+39-44) = -5(-10) \\
  & = 50.
\end{align*}  
\end{solution}

\problem % Created for ILA
Use either the cofactor expansion or row-reduction to find $\det(A)$ for
\[A = \mat{7 3 1 4 2; 0 1 3 4 1; 0 2 10 8 -4; 0 0 0 3 1; 0 0 0 0 7}.\]
\begin{solution}
  The matrix $A$ is almost upper-triangular already, so it may take very little
  row reduction to obtain an upper-triangular matrix.  This gives us the hint
  that row-reduction may be the faster way. \newline \newline %
  The most natural cofactor expansion would be along the first column or last
  row, since each of those has only one nonzero term.  
  \begin{align*}
    \det(A) & = 7 \det\mat{1 3 4 1; 2 10 8 -4; 0 0 3 1; 0 0 0 7}
              = 7 \cdot 7(-1)^{4+4} \det\mat{1 3 4; 2 10 8; 0 0 3} \\
             &  = 49 \det\mat{1 3 4; 2 10 8; 0 0 3} 
             = 49(3)(-1)^{3+3} \det\mat{1 3; 2 10}
              = 147(10-6) = 588.
  \end{align*}
  Row-reduction only requires one step.
  \[A = 
  \mat{7 3 1 4 2; 0 1 3 4 1; 0 2 10 8 -4; 0 0 0 3 1; 0 0 0 0 7}
  \xrightarrow[]{R_3 = R_3 - 2R_2}
  \mat{7 3 1 4 2; 0 1 3 4 1; 0 0 4 0 -6; 0 0 0 3 1; 0 0 0 0 7}.
  \]
  With just one row-replacement (which does not affect the determinant),
  we obtained an upper-triangular matrix, so
  \[\det(A) = 7(1)(4)(3)(7) = 588.\]


\end{solution}

\problem % Created for ILA
A problem from section 2.2, revisited: Find all values of $b$ (if
there are any) so that $\vec{b -3 2}$ is a linear combination of
$\vec{3 1 2}$ and $\vec{1 0 4}$.

\begin{solution}
  Since $\vec{3 1 2}$ and $\vec{1 0 4}$ are linearly independent, we
  know from the Increasing Span Criterion that
  $\left\{ \vec{3 1 2}, \vec{1 0 4}, \vec{b -3 2}\right\}$ will span
  ${\bf R}^3$ (in other words, the three vectors will be linearly
  independent) unless $\vec{b -3 2}$ is in the span of $\vec{3 1 2}$
  and $\vec{1 0 4}$.  We form the matrix whose three columns are the
  vectors and set its determinant equal to $0$.
  \[\det\mat{3 1 b; 1 0 -3; 2 4 2} = 0,\]
  \[3\big[0(2) - (-3)(4)\big] \ - \ 1\big[1(2)-(-3)(2)\big] \ + \
    b\big[1(4)-2(0)\big] = 0,\]
  \[3(12) - 1(8) + b(4) = 0, \qquad 28+4b = 0, \qquad b = -7.\]
\end{solution}

\problem % Dan's notes
For what values of $c$ does the following matrix fail to be invertible?
\[A = \mat{c 1 1; 1 c 1; 1 1 c}.\]
\begin{solution}
We need $\det(A) = 0$.
\[0 = \det(A) = c(c^2-1) - 1(c-1) +1(1-c) = c^3-3c+2 = (c-1)^2(c+2).\] Thus
$c=1$ and $c = -2$ make $A$ fail to be invertible. \newline %
Note: to factor $c^3-3c+2$ above, we did long division after observing that
$c=1$ was a zero of $c^3-3c+2$.  However, we could have instead grouped terms
using the common factor of $c-1$:
\begin{align*}
  c(c^2-1)-1(c-1)+1(1-c)
  & = c(c-1)(c+1)-1(c-1)+1(1-c) \\
  & = (c-1)\Big[c(c+1)-1-1\Big] \\
  & = (c-1)\Big[c^2+c-2\Big] \\
  & = (c-1)(c-1)(c+2).
\end{align*}
\end{solution}

\end{problems} \setcounter{problem}{0}

\newpage

\subsection{4.3, Determinants and volumes}

\begin{problems}
%
\problem % Created for ILA
Suppose we are given two vectors $u$ and $v$ in ${\bf R}^{2}$ and asked to compute
the area of their associated parallelogram, but the computation yields $0$ as the answer.
What does this mean?

\begin{solution}
  If $0 = |\det\mat{u v}|$, then $u$ and $v$ are linearly dependent, so they do
  not actually form a parallelogram.  Instead, since they are on the same line
  through the origin, they only form a portion of a line, which has $0$ area in
  ${\bf R}^{2}$.
\end{solution}

\problem % Created for ILA
  Find the area of the parallelogram determined by the vectors %
  \[v_1 = \vec{1 -3}, \quad v_2 = \vec{2 4}.\]

\begin{solution}
  The area is $\left|\det\mat{v_1 v_2} \right|=
  \left|\det\mat{1 2; -3 4}\right| = \left|1(4)-2(-3)\right| = 10$.
\end{solution}  


\problem % Created for ILA
  Find the area of the parallelogram with vertices
  \[(1, 2), \quad (3, 7), \quad (4, -1),  \quad (6, 4).\]

\begin{solution}
  To use the formula from section 4.2, we need to find vectors $v_1$ and $v_2$
  that determine the area of the parallelogram. For this, we can start at any
  vertex (we will use the first vertex) and find the vectors to each of its two
  adjacent vertices within the parallelogram.  \newline %
  The vector from $(1,2)$ to $(3, 7)$ is $v_1 = \vec{2 5}$. \newline %
  The vector from $(1,2)$ to $(4, -1)$ is $v_2 = \vec{3 -3}$. \newline %
  The area is
  \[\left|\det\mat{v_1 v_2}\right| = \left|\det\mat{2 3; 5 -3}\right| = %
    |2(-3)-3(5)| = |-21| = 21.\]
\end{solution}  


\problem % Created for ILA
  Find the area of the triangle with vertices
  \[(4,-5), \quad (1, 2), \quad (-10, -3).\]

\begin{solution}
  The area is half of the area of the corresponding parallelogram.  To use the
  formula from section 4.2, we need to find vectors $v_1$ and $v_2$ that
  determine the triangle. For this, we can start at any vertex
  (we will use the first vertex) and find the vectors to each of its two
  adjacent vertices within the triangle.  \newline %
  The vector from $(4, -5)$ to $(1, 2)$ is $v_1 = \vec{-3 7}$. \newline %
  The vector from $(4, -5)$ to $(-10, -3)$ is $v_2 = \vec{-14 2}$. \newline %
  The area is
  \begin{align*}
    \dfrac{1}{2} \left|\det\mat{v_1 v_2}\right| %
    & = \dfrac{1}{2} \left|\det\mat{-3 -14; 7 2}\right| %
    = \dfrac{1}{2}\left|(-3)(2)-(-14)(7)\right| \\ %
    & = \dfrac{1}{2}\left|-6 + 98\right| = \dfrac{1}{2}(92) = 46.
  \end{align*}
\end{solution}  


\problem % Created for ILA
Find the volume of the parallelepiped $P$ naturally determined by the three
vectors below.
\[v_1 = \vec{1 0 0}, \quad v_2 = \vec{-1 7 3}, \quad v_3 = \vec{2 4 6}.\]
\begin{solution}
The volume is
\[\left|\det\mat{1 -1 2; 0 7 4; 0 3 6}\right| =
  |1(7(6)-3(4)) - (-1)(0-0) + 2(0-0)| = 30.\]
\end{solution}


\problem % Created for ILA
Find the volume of the parallelepiped $P$ naturally determined by the four
vectors below.
\[v_1 = \vec{4 4 1 2}, \quad v_2 = \vec{0 20 1 0}, \quad v_3 = \vec{-1 -1 1 1},
  \quad v_4 = \vec{5 5 3 4}.\]
\begin{solution}
The volume is
\begin{align*}
  \left|\det\mat{4 0 -1 5; 4 20 -1 5; 1 1 1 3; 2 0 1 4}\right|
  & = \left| 0C_{21} + 20C_{22} + 1C_{32} + 0 C_{42}\right| \\
  & = \left| 20 \det\mat{4 -1 5; 1 1 3; 2 1 4} +
    1(-1)^{3+2} \det \mat{4 -1 5; 4 -1 5; 2 0 4}\right| \\
  & = \left| 20 \Big[4(4-3)+1(4-6)+5(1-2)\Big] -1 (0)\right| = \left|20(4-2-5)\right|\\
  & = 60.
\end{align*}
Above, we did not need to do any work to compute $C_{42}$ because its
$3 \times 3$ matrix had two identical rows and thus had determinant $0$.
\end{solution}


\problem % Created for ILA
Find all values of $h$ so that the volume of the parallelepiped $P$ formed by the
four vectors below is equal to $15$.
\[v_1 = \vec{-2 0 h 0}, \qquad %
  v_2 = \vec{4 1 1 0}, \qquad %
  v_3 = \vec{0 2 1 9}, \qquad %
  v_4 = \vec{6 4 5 1}.\]

\begin{solution}
  There are two values of $h$ that work: $h = 5/82$ and $h = 35/82$.  The volume
  of $P$ is the absolute value of the determinant of the matrix with columns
  $v_1$ through $v_4$, so we set it equal to $15$ and solve for $h$.  We use the
  cofactor expansion along the first column.
  \begin{align*}
    15 & = \left|\det\mat{-2 4 0 6; 0 1 2 4; h 1 1 5; 0 0 9 1}\right| 
         = \left|-2C_{11} + 0 C_{21} + h C_{31} + 0 C_{41}\right| \\
       & = \left|-2(-1)^{1+1} \det\mat{1 2 4; 1 1 5; 0 9 1} + h(-1)^{3+1}
         \det\mat{4 0 6; 1 2 4; 0 9 1}\right| \\
       & = \left| -2 \Big[1(1-9\cdot 5)-2(1-0)+4(9-0)\Big]
         + h \Big[4(2\cdot 1-9 \cdot 4)-0+6(1 \cdot 9 -0)\Big] \right| \\
       & = \left| -2(-44-2+36) + h(-136+54)\right| \\
       & = \left|-2(-10) - 82h\right| = \left|20-82h\right|.         
  \end{align*}
  Therefore, $\left|20-82h\right| = 15$. \newline
  Case 1: $20-82h = 15$, so $82h = 5$ and $h = 5/82$. \newline %
  Case 2: $-(20-82h) = 15$, so $-20+82h = 15$, thus $82h = 35$
  and $h = 35/82$.
\end{solution}


\problem % Created for ILA
Let $R$ be the triangle with vertices $(0, 0)$, $(1, 4)$, and $(10, 20)$.  Find
the area of the new triangle $T(R)$, where $T: {\bf R}^2 \to {\bf R}^2$ is the
transformation $T(x) = \mat{8 -3; 1 2}$.

\begin{solution}
  The area of $R$ is
  \[\dfrac{1}{2} \left| \det\mat{1 10; 4 20}\right| = \dfrac{1}{2}|-20| = 10.\]
  Let $A = \mat{8 -3; 1 2}$, noting $|\det(A)| = 19$.  From the area formula from
  chapter 4, we know
  \[\textup{Area}(T(R)) = \left|\det(A)\right|\cdot \textup{Area}(R)
    = 19(10) = 190.\]
  Another way to do this problem is to transform the vertices of the original triangle
  into the new triangle: we compute $T\vec{0 0} = \vec{0 0}$, 
  \[T\vec{1 4} = \mat{8 -3; 1 2} \vec{1 4} = \vec{-4 9}, \qquad
    T\vec{10 20} = \mat{8 -3; 1 2} \vec{10 20} = \vec{20 50}.\]
  Therefore, the vertices of $T(R)$ are $(0, 0)$, $(-4,9)$, and $(20,50)$.
  \[\textup{Area}(T(R)) = \dfrac{1}{2}
    \left|\det\mat{-4 20; 9 50}\right| = \dfrac{1}{2} |-200-180| = \dfrac{1}{2}(380)
    = 190.\]
\end{solution}


\problem % Created for ILA
Let $T: {\bf R}^2 \to {\bf R}^2$ be the transformation
\[T(x) = \mat{7 -1; 0 3} x.\] Find the area of $T(C)$, where $C$ is the circle
of radius $2$ centered at the origin.

\begin{solution}
With $A = \mat{7 -1; 0 3}$, we use the area formula from section 4.3:
\[\textup{Area}(T(C)) = \left| \det(A) \right| \textup{Area}(C)
  = 21 \cdot (\pi(2^2)) = 84 \pi.\]
\end{solution}


\problem % Dan's notes
  If the area of a fish (in a photo) is $7$ square inches, and if we apply the
  shear transformation $T(x) = \mat{1 13; 0 1}$ to the photo, what is the area
  of the fish in the new photo?

\begin{solution}
  To get the area of the new fish, we multiply the original area by the
  determinant of the standard matrix for $T$.  The new area is
  $\det\mat{1 13; 0 1} \cdot 7 = 1 \cdot 7 = 7$.
\end{solution}


%\problem % Dan's notes
%Find the volume of the triangular pyramid with the vertices 
%\[(0,0,0), \quad (0,0,1), \quad (1,0,0), \quad (1,2,3).\]
%
%\begin{solution}
%  The volume of the triangular pyramid is half the volume of the corresponding
%  rectangular pyramid.  We use the vectors from $(0,0,0)$ to each of the three
%  vertices in order to create the columns of the corresponding matrix.  The
%  volume is
%  \[\dfrac{1}{2} \det\mat{0 1 1; 0 0 2; 1 0 3} =
%    \dfrac{1}{2} \left(0 - 1(0-2) + 1(0-0)\right) = 1.\]
%\end{solution}

\problem % Made for ILA
Suppose $u$ and $v$ are vectors in ${\bf R}^2$ that naturally form a
parallelogram with area $10$ square units.  What is the area of the
parallelogram naturally formed by $3u$ and $-v$. Justify this answer
in two ways: first by using geometry, then by using properties of
determinants.

\begin{solution}
  Geometrically: stretching $u$ by a factor of three will scale the
  area by a factor of $3$, whereas multiplying $v$ by $-1$ will only
  change the direction of facing of the parallelogram.  Therefore, the
  resulting area is $3(10)=30$.
  % 
  Determinants: the area of the parallelogram determined by $u$ and
  $v$ is
   \[\textup{Area} = \left|\det \mat{- u -; - v -}\right|=10 \textup{
       units}^2.\] Scaling a row of a matrix by $3$ multiplies the
   determinant by $3$, and multiplying a row by $-1$ changes the
   determinant by $-1$ (this will be irrelevant since we take the
   absolute value).
   \[\textup{Area} = \left|\det \mat{- 3u -; - -v -}\right| =
     |10(3)(-1)| = 30 \textup{ units}^2.\]
\end{solution}

 
 
\end{problems}\setcounter{problem}{0}
\newpage